\documentclass[12pt,a4paper]{article}
\usepackage[utf8]{inputenc}
\usepackage[indonesian]{babel}
\usepackage{amsmath}
\usepackage{amsfonts}
\usepackage{amssymb}
\usepackage{geometry}

\geometry{margin=2.5cm}

\title{Menentukan Persamaan Kartesius dari Persamaan Parametrik}
\author{Ahmad Fakhrul Bawani}
\date{}

\begin{document}

\maketitle

\section*{Soal}

The parametric equations and parameter intervals for the motion of a particle in the $xy$-plane are given below. Identify the particle's path by finding a Cartesian equation for it. Graph the Cartesian equation. Indicate the portion of the graph traced by the particle and the direction of motion.
\begin{equation*}
  x = 5\cos(2t), \quad y = 5\sin(2t), \quad 0 \le t \le \frac{\pi}{2}
\end{equation*}

\subsection*{1. Mencari Persamaan Kartesius}
Kita dapat mengeliminasi parameter $t$ dengan mengisolasi fungsi trigonometri pada masing-masing persamaan:
\begin{align*}
  x = 5\cos(2t) &\implies \cos(2t) = \frac{x}{5} \\
  y = 5\sin(2t) &\implies \sin(2t) = \frac{y}{5}
\end{align*}

Gunakan identitas trigonometri utama $\cos^2(\theta) + \sin^2(\theta) = 1$ di mana $\theta = 2t$:
\begin{align*}
  \cos^2(2t) + \sin^2(2t) &= 1 \\
  \left(\frac{x}{5}\right)^2 + \left(\frac{y}{5}\right)^2 &= 1 \\
  \frac{x^2}{25} + \frac{y^2}{25} &= 1 \\
  x^2 + y^2 &= 25
\end{align*}

Jadi, persamaan Kartesius untuk jalur partikel tersebut adalah \textbf{$x^2 + y^2 = 25$}, yang merepresentasikan sebuah lingkaran yang berpusat di titik asal $(0,0)$ dengan jari-jari $R = 5$.

\subsection*{2. Menentukan Porsi Grafik dan Arah Pergerakan}
Untuk mengetahui bagian lingkaran mana yang dilalui dan ke mana arahnya, kita substitusikan nilai batas parameter $t$ ke dalam persamaan parametrik awal:

\begin{itemize}
  \item \textbf{Saat Titik Awal ($t = 0$):}
    \begin{align*}
      x &= 5\cos(2 \cdot 0) = 5\cos(0) = 5(1) = 5 \\
      y &= 5\sin(2 \cdot 0) = 5\sin(0) = 5(0) = 0
    \end{align*}
    Partikel mulai bergerak dari titik \textbf{$(5, 0)$}.

  \item \textbf{Saat Titik Tengah ($t = \frac{\pi}{4}$):}
    \begin{align*}
      x &= 5\cos\left(2 \cdot \frac{\pi}{4}\right) = 5\cos\left(\frac{\pi}{2}\right) = 5(0) = 0 \\
      y &= 5\sin\left(2 \cdot \frac{\pi}{4}\right) = 5\sin\left(\frac{\pi}{2}\right) = 5(1) = 5
    \end{align*}
    Partikel berada di titik \textbf{$(0, 5)$}.

  \item \textbf{Saat Titik Akhir ($t = \frac{\pi}{2}$):}
    \begin{align*}
      x &= 5\cos\left(2 \cdot \frac{\pi}{2}\right) = 5\cos(\pi) = 5(-1) = -5 \\
      y &= 5\sin\left(2 \cdot \frac{\pi}{2}\right) = 5\sin(\pi) = 5(0) = 0
    \end{align*}
    Partikel berhenti di titik \textbf{$(-5, 0)$}.
\end{itemize}

\subsection*{Kesimpulan Analisis}
Berdasarkan uji titik koordinat di atas, karakteristik pergerakan partikel dapat disimpulkan sebagai berikut:
\begin{itemize}
  \item \textbf{Persamaan Kartesius:} $x^2 + y^2 = 25$
  \item \textbf{Porsi Grafik yang Dilalui:} Hanya bagian \textbf{setengah lingkaran atas} (berada di kuadran I dan kuadran II), karena nilai koordinat bergerak dari $x = 5$ menuju $x = -5$ dengan nilai $y \ge 0$.
  \item \textbf{Arah Gerakan:} Bergerak berlawanan arah jarum jam (\textbf{\textit{counterclockwise}}), dimulai dari titik kanan $(5,0)$ melewati puncak atas $(0,5)$ hingga berakhir di titik kiri $(-5,0)$.
\end{itemize}

\end{document}